\chapter{Reference: commands, files, exit codes}\label{ch:reference}

\section{Every \texttt{ate} command}

\begingroup\footnotesize
\begin{longtable}{@{}>{\raggedright\arraybackslash}p{5.2cm}>{\raggedright\arraybackslash}p{9.6cm}@{}}
\toprule
Command & What it does \\
\midrule
\endhead
\texttt{ate} & Opens the application. With an unusable config or database, refuses and exits 1. \\
\texttt{ate --config PATH} & Read that config instead of the platform default. It may sit before the
  subcommand (\texttt{ate --config X db status}) or after it (\texttt{ate db status --config X}). \\
\texttt{ate -V}, \texttt{ate --version} & Print the version and exit 0. \\
\texttt{ate --help} & Print usage at any level and exit 0. \\
\texttt{ate --allow-unsandboxed-pdf} & When ATE is root or an administrator, render Test Report PDFs
  without Chromium's sandbox instead of asking (or skipping them). Has no effect for an ordinary
  user, and says so. \\
\texttt{ate db status [PATH]} & Print the database's schema version. Read-only: it never creates a
  database. Exits 1 with the refusal if it is missing, outdated, unreadable or not a database. \\
\texttt{ate db migrate [PATH] [--yes|-y]} & Create or upgrade the database, after confirmation unless
  \texttt{--yes}. Copies the database beside itself first, when a migration will run. Refuses a
  database built by a newer ATE, and refuses to migrate when the copy cannot be written. \\
\texttt{ate setup} & Create the Station's directories, write a default config if none exists, and
  migrate the database named by the config. \\
\texttt{--root DIR} & Prepare DIR as the Station root: DIR/data and DIR/config. \\
\texttt{--data-dir DIR} & Use DIR for the data directory instead of the root's \texttt{data/}. \\
\texttt{--config-dir DIR} & Use DIR for the config directory instead of the root's \texttt{config/}. \\
\texttt{--no-permissions} & Skip the group, ownership and ACL treatment explicitly (packagers and
  test runs). \\
\texttt{ate definition update <file.atedef>} & Recompute every listed file's SHA-256 in place,
  filling in missing ones. Rewrites only if something changed. \\
\texttt{ate \_pdf <html> <pdf> --paper A4|Letter [--unsandboxed]} & Hidden helper ATE starts as a
  child process to render one Test Report. Not for humans to run. \\
\texttt{ate \_selftest} & Hidden: the frozen bundle inventories what it carries and exits 1 if
  anything is missing. \\
\texttt{ate \_selftest --run} & Hidden: runs a built-in one-step Test Schedule head end to end in a
  temporary directory, and checks both reports were written. A privileged run needs
  \texttt{--allow-unsandboxed-pdf} to produce the PDF. \\
\texttt{ate \_selftest --scan DIR} & Hidden: release CI's headless launch. Reads the config
  \texttt{ate setup} wrote, refuses a database at the wrong schema version, and reports the schedules
  in a Test Definition directory and the broken ones. Writes nothing. \\
\bottomrule
\end{longtable}
\endgroup

The hidden commands are deliberately not in \texttt{ate --help}; they are what ATE and the packager
call, not what an Operator types. An administrator can still run them, and each one prints what it
found and exits 1 if anything is missing or broken. \texttt{\_selftest} inventories the bundle,
\texttt{\_selftest --run} runs a built-in one-step Test Schedule in its own temporary directory,
\texttt{\_selftest --scan DIR} checks the database version and reports the Test Definitions in a
directory, and \texttt{\_pdf} is the child process the report writer starts:

\begin{lstlisting}[language=ATEShell]
ate --help
ate -V
ate _selftest
ate _selftest --run
ate _selftest --scan definitions/MR-1043
ate _pdf report.html report.pdf --paper A4
\end{lstlisting}

\section{Where things are}

\begin{center}
\footnotesize
\begin{tabular}{@{}>{\raggedright\arraybackslash}p{4.0cm}>{\raggedright\arraybackslash}p{5.4cm}>{\raggedright\arraybackslash}p{6.0cm}@{}}
\toprule
Thing & Windows & Linux \\
\midrule
Config & \texttt{\%PROGRAMDATA\%\textbackslash ATE}\allowbreak\texttt{\textbackslash config.json} & \texttt{/etc/ate/config.json} \\
Data directory (\texttt{<data>}) & \texttt{\%PROGRAMDATA\%\textbackslash ATE} & \texttt{/var/lib/ate} \\
Database & \texttt{<data>/ate.db} & \texttt{/var/lib/ate/ate.db} \\
Test Reports & \texttt{<data>/reports/}\allowbreak\texttt{<UUT Type>} & \texttt{/var/lib/ate/reports/}\allowbreak\texttt{<UUT Type>} \\
Application log & \texttt{<data>/logs} & \texttt{/var/lib/ate/logs} \\
Media Releases & \texttt{<data>/definitions} & \texttt{/var/lib/ate/definitions} \\
Program & \texttt{C:\textbackslash Program Files\textbackslash ATE} & \texttt{/opt/ate}, linked as \texttt{/usr/bin/ate} \\
Licence notices & \texttt{<program}\textbackslash\texttt{\_internal}\allowbreak\texttt{\textbackslash licences} & \texttt{<program/\_internal/}\allowbreak\texttt{licences} \\
Run log & in the database, table \texttt{log\_line} & the same \\
\bottomrule
\end{tabular}
\end{center}

The log and report paths are what the config says; the defaults above are what \texttt{ate setup}
writes, and it writes them already resolved to absolute paths. A path left relative in the config is
resolved against the config file's directory, so only a hand-written config with relative paths keeps
its data next to the config.

\section{What this manual cannot tell you}

Be honest about these, because a Station's problem is often in one of them:

\begin{itemize}
  \item \textbf{What the interface looks like, and what a click does.} ATE has never been run with a
    display that a human watched. Layout, fonts, dialog appearance and the exact wording of prompts
    other than the ones quoted here are unverified, and the appearance work is still in progress.
  \item \textbf{What Windows makes of an unsigned installer.} There is no signing certificate, so the
    artifacts are unsigned. The installer installed, upgraded and uninstalled on a CI runner with
    \texttt{/VERYSILENT}, which suppresses dialogs; what a person sees when they double-click it is
    not known here.
  \item \textbf{Real instruments.} Every Device test in the suite runs against a loopback peer, a
    \texttt{socket://} URL standing in for a serial port, or pyvisa-py's own backends. Nothing has
    been connected to GPIB, USBTMC or RS-232 hardware, and NI-VISA is not shipped. Timing,
    terminator behaviour and error handling on a real instrument are untested.
  \item \textbf{Why the one Windows stall happened.} The fix is tested; the cause is a hypothesis
    about the runner's disk.
  \item \textbf{Whether the wheel carries the licence notices.} The bundle, the tarball and the
    \texttt{.deb} demonstrably do. The wheel is not checked for them.
\end{itemize}

\cando{look up any command, any file, and any claim this manual does not actually support.}
